\documentclass[a4paper,12pt]{article}
\usepackage[utf8]{inputenc}
\usepackage[T1]{fontenc}
\usepackage[magyar]{babel}
\usepackage{amsmath,amssymb,amsthm}
\usepackage[margin=2.5cm]{geometry}
\usepackage{parskip}
\usepackage{enumitem}
\usepackage{xcolor}
\usepackage{mdframed}
\usepackage{booktabs}
\usepackage{hyperref}
\hypersetup{colorlinks=true,linkcolor=blue}

\definecolor{taskbg}{rgb}{0.94,0.97,1.0}
\definecolor{darkgray}{rgb}{0.4,0.4,0.4}

\title{\LARGE\bfseries Numerikus módszerek\\[6pt]
  \large 9.\ hét -- Papíron számolós feladatok}
\author{}
\date{2025/26.\ tanév}

\begin{document}
%\maketitle
%\tableofcontents
%\bigskip\hrule\bigskip

\section{Gauss--Seidel és SOR iteráció}

\begin{mdframed}[backgroundcolor=taskbg,linewidth=0pt,innertopmargin=8pt,innerbottommargin=8pt]
\textbf{Optimális SOR-tétel (Young-tétel).} \quad
{\normalfont\small SOR = \textit{Successive Over-Relaxation} (egymást követő túlrelaxáció)}

\medskip
\textit{Feltételek:}
\begin{enumerate}[noitemsep,topsep=2pt]
  \item $A$ szimmetrikus, pozitív definit;
  \item $A$ tridiagonális szerkezetű;
  \item a Jacobi-iteráció $B_J$ mátrixának sajátértékei valósak.
\end{enumerate}

\medskip
\textit{Állítás:}\ Ha $\varrho(B_J)<1$, akkor
\[
  \omega_0 = \frac{2}{1+\sqrt{1-\varrho(B_J)^2}}
\]
az optimális relaxációs paraméter, és a hozzá tartozó SOR-iteráció spektrálsugara
\[
  \varrho(B_{\omega_0}) = \omega_0 - 1.
\]
Bármely $\omega\in(0,2)$-re $\varrho(B_\omega)<1$ (konvergencia), és
$\varrho(B_{\omega_0})\le\varrho(B_\omega)$ (optimalitás).
\end{mdframed}

\subsection*{1. feladat -- Gauss--Seidel lépések kézzel}

Legyen
$$A = \begin{bmatrix}4&-1\\-1&4\end{bmatrix}, \quad b = \begin{bmatrix}3\\7\end{bmatrix},
\quad x^{(0)} = \begin{bmatrix}0\\0\end{bmatrix}.$$

\textbf{Végezz 2 Gauss--Seidel-iterációs lépést kézzel!}

%\textit{Emlékeztető:} Az $i$-edik egyenletből:
%$$x_i^{(k+1)} = \frac{1}{a_{ii}}\left(b_i - \sum_{j<i}a_{ij}x_j^{(k+1)} - \sum_{j>i}a_{ij}x_j^{(k)}\right)$$


\subsection*{2. feladat -- Optimális SOR paraméter meghatározása}

Legyen
$$A = \begin{bmatrix}2&-1&0\\-1&2&-1\\0&-1&2\end{bmatrix}.$$

\begin{enumerate}[noitemsep,topsep=2pt]
  \item Állapítsd meg, hogy alkalmazható-e az optimális SOR-tétel!
  \item Számítsd ki $\varrho(B_J)$-t (a Jacobi-iteráció spektrálsugarát)!
  \item Add meg az optimális $\omega_0$-t!
\end{enumerate}


\subsection*{3. feladat -- SOR-tétel alkalmazhatósága}

Döntsd el az alábbi mátrixokról, hogy alkalmazható-e rájuk az optimális SOR-tétel!
Ha igen, add meg $\omega_0$-t.

(a) $A_1 = \begin{bmatrix}3&-1&0\\-1&3&-1\\0&-1&3\end{bmatrix}$

(b) $A_2 = \begin{bmatrix}2&1\\-1&3\end{bmatrix}$

(c) $A_3 = \begin{bmatrix}4&-1&-1&0\\-1&4&0&-1\\-1&0&4&-1\\0&-1&-1&4\end{bmatrix}$


\section{Richardson-iteráció}
\subsection*{4. feladat -- Richardson konvergencia tartomány}

$$A = \begin{bmatrix}5&1\\1&5\end{bmatrix}, \quad b = \begin{bmatrix}6\\6\end{bmatrix}$$

Add meg azt a $p\in\mathbb{R}$ értéktartományt, amelyre a Richardson-iteráció $R(p)$ konvergens!


\subsection*{5. feladat -- Optimális Richardson paraméter}

Az előző feladat mátrixára ($A=\begin{bmatrix}5&1\\1&5\end{bmatrix}$):

\begin{enumerate}[noitemsep,topsep=2pt]
  \item Add meg az optimális $p_0$-t!
  \item Add meg a hozzá tartozó $q$ kontrakciós együtthatót!
  \item Hány lépés szükséges $10^{-4}$ pontossághoz, ha $\|x^{(0)}-x^*\|\le 1$?
\end{enumerate}


\subsection*{6. feladat -- Richardson-lépések kézzel}

$$A = \begin{bmatrix}5&1\\1&5\end{bmatrix},\quad b = \begin{bmatrix}6\\6\end{bmatrix},
\quad x^{(0)} = \begin{bmatrix}0\\0\end{bmatrix},\quad p = \frac{1}{5}$$

Végezz 2 Richardson-lépést kézzel (reziduum-vektoros alakban)!


\section{Bolzano-tétel és intervallumfelezés}
\subsection*{7. feladat -- Intervallumfelezés kézzel}

$$f(x) = x^3 - 2x - 5$$, intervallum: $$[2, 3]$$.

Végezz \textbf{3 intervallumfelezési lépést} kézzel! Töltsd ki a táblázatot:

\begin{center}
\begin{tabular}{l|l|l|l|l|l}
\hline
k & a & b & c=(a+b)/2 & f(c) & új intervallum\\
\hline
0 & 2 & 3 & 2.5 & -0.375 & [2.5, 3]\\
1 & 2.5 & 3 & ... & ... & ...\\
2 & ... & ... & ... & ... & ...\\
\hline
\end{tabular}
\end{center}

%\textit{Ellenőrzés:} $f(2)=-1<0$, $f(3)=16>0$ \checkmark

\subsection*{8. feladat -- Szükséges lépésszám}

Hány intervallumfelezési lépés szükséges, hogy az $[1, 2]$ intervallumon $10^{-4}$ pontosságot
elérjünk?


\subsection*{9. feladat -- Bolzano-tétel alkalmazása}

Állapítsd meg, van-e gyöke az $f(x) = e^x - 3x^2$ függvénynek a $[0,1]$ intervallumon!
Alkalmazható-e a Bolzano-tétel?


\section{Fixponttételek és egyszerő iterációk}
\subsection*{10. feladat -- Kontrakció igazolása}

Igazold, hogy $\varphi(x) = \dfrac{\cos(x)}{2} + 1$ kontrakció a $[0, 2]$ intervallumon!
Add meg a kontrakciós együtthatót $q$-t!


\subsection*{11. feladat -- Banach-tétel feltételei}

$\varphi(x) = \dfrac{x^2}{3} + \dfrac{1}{3}$

\begin{enumerate}[noitemsep,topsep=2pt]
  \item Leképezi-e $[0,1]$-t önmagára?
  \item Kontrakció-e?
  \item Add meg a fixpontot!
  \item Hány lépés kell 0.01-es pontossághoz $x_0 = 0.5$-ből indulva?
\end{enumerate}


\section{Konvergencia rend}
\subsection*{12. feladat -- Konvergencia rend meghatározása}

Határozd meg az alábbi nullsorozat konvergencia rendjét!

$$x_k = \frac{1}{3^k}$$


\subsection*{13. feladat -- Konvergencia rend Taylor-sorral}

Az $x^3-x-1=0$ egyenlet $\varphi(x)=\sqrt[3]{x+1}$ iterációjának konvergencia rendje?


\section{Newton-, húr- és szelőmódszer}
\subsection*{14. feladat -- Newton-lépések kézzel}

$f(x) = x^3 - 2$, $x_0 = 2$.

Végezz \textbf{2 Newton-lépést} kézzel!


\subsection*{15. feladat -- Húrmódszer kézzel}

$f(x) = x^2 - 3$, $[a,b]=[1,2]$.

Végezz \textbf{2 húrmódszer-lépést} kézzel!


\subsection*{16. feladat -- Szelőmódszer kézzel}

$f(x) = x^2 - 3$, $x_0=1$, $x_1=2$.

Végezz \textbf{2 szelőmódszer-lépést} kézzel!


\end{document}
